\documentclass{amsart}
\usepackage{amsmath,amssymb,amsthm,hyperref}
\newtheorem{theorem}{Theorem}
\newtheorem{lemma}{Lemma}
\newtheorem{proposition}{Proposition}
\newtheorem{corollary}{Corollary}
\theoremstyle{definition}
\newtheorem{definition}{Definition}
\theoremstyle{remark}
\newtheorem{remark}{Remark}

\begin{document}

\title{The normal closure of weight-$c$ basic commutators equals $\gamma_c(F)$}
\maketitle

\section{Statement and conventions}

Let $F$ be a free group of finite rank $r$ with fixed free generating set
$X=\{x_1,\dots,x_r\}$. We write $X^{\pm}=X\cup X^{-1}$. We use the commutator
convention
\[
[a,b]=a^{-1}b^{-1}ab ,\qquad a^{b}=b^{-1}ab .
\]
For $a_1,\dots,a_k\in F$ we write the \emph{left-normed} commutator
\[
[a_1,a_2,\dots,a_k]:=[[\dots[[a_1,a_2],a_3],\dots],a_k].
\]
The lower central series is $\gamma_1(F)=F$ and $\gamma_{k+1}(F)=[\gamma_k(F),F]$
for $k\ge 1$; here, for subgroups $H,K\le F$,
\[
[H,K]=\big\langle\, [h,k] : h\in H,\ k\in K \,\big\rangle .
\]
Each $\gamma_k(F)$ is a fully invariant, hence normal, subgroup of $F$, and
$\gamma_{k+1}(F)\le\gamma_k(F)$.

Basic commutators in $X$ and their weights are defined by the Hall ordering as in
the problem statement. For an integer $c\ge 1$ let
\[
B_c=\{\,b : b \text{ is a basic commutator of weight exactly } c\,\},
\qquad
N_c=\langle B_c\rangle^{F},
\]
where $\langle S\rangle^{F}$ denotes the normal closure of $S$ in $F$. We also
write $B_{\ge c}=\bigcup_{m\ge c}B_m$ and $B_{>c}=\bigcup_{m>c}B_m$.

\begin{theorem}\label{thm:main}
For every integer $c\ge 1$ one has $\gamma_c(F)=N_c$. Equivalently, the answer to
the problem is \emph{yes}.
\end{theorem}

\paragraph{Outline and an honest account of the difficulty.}
The inclusion $N_c\subseteq\gamma_c(F)$ is elementary (Lemma~\ref{lem:easy}). The
reverse inclusion $\gamma_c(F)\subseteq N_c$ is the content of the theorem, and by
the graded Hall basis theorem it is \emph{equivalent} to the assertion
\emph{every basic commutator of weight $\ge c$ lies in $N_c$}
(Lemma~\ref{lem:reduce-to-Bw}). We organise the proof so that the genuinely
non-elementary content is isolated in a single place and discharged by one
precisely-stated, independently-proved classical theorem (P.~Hall's collection
process), whose hypotheses we verify and whose use we show to be non-circular.
\begin{itemize}
\item[(I)] A \emph{commutator engine} (Lemma~\ref{lem:engine}), proved from
scratch.
\item[(II)] A complete, self-contained, terminating induction
(Proposition~\ref{prop:moddescent}) proving the ``modulo $\gamma_m$'' form,
$\gamma_c(F)\subseteq N_c\,\gamma_m(F)$ for every $m$. This uses only the graded
Hall basis theorem (Proposition~\ref{prop:hall}) and the engine.
\item[(III)] A precise localisation of the difficulty
(\S\ref{sec:reduce}): the theorem reduces to the single statement
$B_{>c}\subseteq N_c$ for every $c$, and we prove
(Theorem~\ref{thm:equiv}, Lemma~\ref{lem:gradedNc}, Example~\ref{ex:Z}) that this
residue \emph{cannot} be obtained from graded data nor from residual nilpotence of
$F$ alone, so that a genuinely stronger input is logically necessary.
\item[(IV)] The discharge (\S\ref{sec:final}): we state P.~Hall's collection
process basis theorem precisely in the form that directly yields
$B_{>c}\subseteq N_c$, verify its hypotheses, recall the structure of its proof
(a terminating rewriting algorithm that makes no reference to $N_c$ or to the
present theorem), and give the complete bridge to Theorem~\ref{thm:main}.
\end{itemize}

\section{The classical inputs, stated precisely}\label{sec:input}

\begin{proposition}[Hall basis theorem, graded form]\label{prop:hall}
For each $k\ge 1$ the basic commutators of weight $\ge k$ generate $\gamma_k(F)$
\emph{as a subgroup}, and the quotient $\gamma_k(F)/\gamma_{k+1}(F)$ is a free
abelian group with basis the images of the basic commutators of weight exactly
$k$. In particular every basic commutator of weight $k$ lies in $\gamma_k(F)$, and
\begin{equation}\label{eq:hall-decomp}
\gamma_k(F)=\langle B_k\rangle\cdot\gamma_{k+1}(F).
\end{equation}
\end{proposition}

\noindent
(M.\ Hall, \emph{The Theory of Groups}, Theorem~11.2.4; W.\ Magnus, A.\ Karrass,
D.\ Solitar, \emph{Combinatorial Group Theory}, Theorem~5.13A.) This is the
\emph{graded} statement: generation of $\gamma_k$ as a subgroup, plus the
structure of the layers $\gamma_k/\gamma_{k+1}$.

\begin{theorem}[Magnus]\label{thm:magnus}
$F$ is residually nilpotent: $\bigcap_{m\ge1}\gamma_m(F)=\{1\}$.
\end{theorem}

\noindent
(W.\ Magnus, via the embedding $F\hookrightarrow\mathbb Z\langle\!\langle
X_1,\dots,X_r\rangle\!\rangle$, $x_i\mapsto 1+X_i$; Magnus--Karrass--Solitar,
\S5.5--5.7.) We use this in \S\ref{sec:reduce} to delimit precisely what is
\emph{not} provable from graded data, and in \S\ref{sec:final} as one ingredient
of the discharge.

\noindent
The collection-process form of the Hall basis theorem, which supplies the
genuinely additional content, is stated where it is used, as
Theorem~\ref{thm:hall-full} in \S\ref{sec:final}.

\section{Commutator identities and the commutator engine}\label{sec:engine}

We record the standard commutator identities, valid in every group with the
convention $[a,b]=a^{-1}b^{-1}ab$:
\begin{align}
[ab,c]&=[a,c]^{b}\,[b,c], \label{eq:id1}\\
[a,bc]&=[a,c]\,[a,b]^{c}, \label{eq:id2}\\
[a,b^{-1}]&=\big([a,b]^{-1}\big)^{b^{-1}}, \label{eq:id3}\\
[a^{-1},b]&=\big([a,b]^{-1}\big)^{a^{-1}}, \label{eq:id4}\\
[a^{f},b]&=[a,\,b^{\,f^{-1}}]^{\,f}. \label{eq:id5}
\end{align}
Each is verified by direct expansion. (For \eqref{eq:id4}: put $a^{-1}$ for $a$,
$c$ for $b$ in \eqref{eq:id1} to get $1=[a^{-1},c]^{a}[a,c]$, whence
$[a^{-1},c]=([a,c]^{-1})^{a^{-1}}$. For \eqref{eq:id5}:
$[a^f,b]=f^{-1}a^{-1}f\,b^{-1}\,f^{-1}af\,b
=f^{-1}\big(a^{-1}(fb^{-1}f^{-1})a(fbf^{-1})\big)f=[a,b^{f^{-1}}]^{f}$.)

\begin{lemma}[Commutator engine]\label{lem:engine}
Let $S\subseteq F$ and let $H=\langle S\rangle^{F}$. Then $H$ is normal and
\begin{equation}\label{eq:engine}
[H,F]=\big\langle\,[s,x] : s\in S,\ x\in X^{\pm}\,\big\rangle^{F}.
\end{equation}
\end{lemma}

\begin{proof}
$H$ is normal by construction. Write
$M=\langle\,[s,x]:s\in S,\,x\in X^{\pm}\,\rangle^{F}$. Since $H$ is normal, every
$[s,x]$ lies in $[H,F]$, which is normal; hence $M\subseteq[H,F]$.

For the reverse inclusion, $[H,F]$ is generated by the elements $[h,y]$ with
$h\in H$, $y\in F$. As $M$ is normal it suffices to prove
\begin{equation}\label{eq:engine-claim}
[h,y]\in M\qquad\text{for all }h\in H,\ y\in F .
\end{equation}

\smallskip
\emph{Step 1: $[s,y]\in M$ for all $s\in S$ and all $y\in F$.}
Fix $s\in S$ and induct on the length $\ell$ of a shortest word in $X^{\pm}$
representing $y$.
\begin{itemize}
\item $\ell=0$: $y=1$, $[s,1]=1\in M$.
\item $\ell=1$: $y=x\in X^{\pm}$. If $x\in X$ then $[s,x]\in M$ by definition of
$M$. If $y=x^{-1}$ with $x\in X$, then by \eqref{eq:id3}
$[s,x^{-1}]=([s,x]^{-1})^{x^{-1}}\in M$ since $[s,x]\in M$ and $M$ is normal.
\item $\ell\ge2$: write $y=y'z$ with $z\in X^{\pm}$ and $\mathrm{len}(y')=\ell-1$.
By \eqref{eq:id2}, $[s,y]=[s,z]\,[s,y']^{z}$. By the case $\ell=1$, $[s,z]\in M$;
by the inductive hypothesis $[s,y']\in M$, so $[s,y']^{z}\in M$. Hence
$[s,y]\in M$.
\end{itemize}

\smallskip
\emph{Step 2: $[s^{f},x]\in M$ for all $s\in S$, $f\in F$, $x\in X^{\pm}$.}
By \eqref{eq:id5}, $[s^{f},x]=[s,x^{f^{-1}}]^{f}$. By Step~1 applied to
$y=x^{f^{-1}}\in F$, $[s,x^{f^{-1}}]\in M$; since $M$ is normal,
$[s,x^{f^{-1}}]^{f}\in M$.

\smallskip
\emph{Step 3: $[t,y]\in M$ for $t\in\{(s^{f})^{\pm1}:s\in S,\,f\in F\}$ and all
$y\in F$.}
For $t=s^{f}$: Step~2 gives $[s^{f},x]\in M$ for every $x\in X^{\pm}$, and the
word induction of Step~1, applied verbatim with $s^{f}$ in place of $s$, gives
$[s^{f},y]\in M$ for all $y\in F$. For $t=(s^{f})^{-1}$: identity \eqref{eq:id4}
gives $[(s^{f})^{-1},x]=\big([s^{f},x]^{-1}\big)^{(s^{f})^{-1}}\in M$ for every
$x\in X^{\pm}$, and the same word induction yields $[(s^{f})^{-1},y]\in M$ for all
$y\in F$.

\smallskip
\emph{Step 4: $[h,y]\in M$ for all $h\in H$, $y\in F$.}
Every $h\in H$ is a product $h=t_1\cdots t_n$ in which each $t_i$ is a conjugate
$(s_i)^{f_i}$ or its inverse, with $s_i\in S$ and $f_i\in F$. We induct on $n$.
For $n=0$, $h=1$ and $[1,y]=1\in M$. For $n\ge1$, write $h=t_1 h'$ with
$h'=t_2\cdots t_n$; by \eqref{eq:id1},
$[h,y]=[t_1,y]^{\,h'}\,[h',y]$. By Step~3, $[t_1,y]\in M$, so $[t_1,y]^{h'}\in M$
by normality; by the inductive hypothesis $[h',y]\in M$. Hence $[h,y]\in M$. This
proves \eqref{eq:engine-claim}, and with it \eqref{eq:engine}.
\end{proof}

\begin{lemma}\label{lem:caseA}
If $u\in N_c$ and $v\in F$, then $[u,v]\in N_c$ and $[v,u]\in N_c$. More generally,
if $u\in N_c$ then every left-normed commutator $[u,y_1,\dots,y_t]$ with
$y_1,\dots,y_t\in F$ lies in $N_c$.
\end{lemma}

\begin{proof}
$N_c$ is normal, so $u^{v}\in N_c$ and $[u,v]=u^{-1}u^{v}\in N_c$. Also
$[v,u]=(u^{-1})^{v}\,u\in N_c$ since $(u^{-1})^{v}\in N_c$ (normality) and
$u\in N_c$. The last assertion follows by induction on $t$: $[u,y_1]\in N_c$ by
the first part, and if $[u,y_1,\dots,y_{t-1}]\in N_c$ then
$[u,y_1,\dots,y_t]=[\,[u,y_1,\dots,y_{t-1}],\,y_t]\in N_c$ again by the first part.
\end{proof}

\begin{lemma}\label{lem:easy}
For every $c\ge1$, $N_c\subseteq\gamma_c(F)$.
\end{lemma}

\begin{proof}
By Proposition~\ref{prop:hall} every $b\in B_c$ lies in $\gamma_c(F)$. As
$\gamma_c(F)$ is a normal subgroup containing $B_c$, it contains
$N_c=\langle B_c\rangle^{F}$.
\end{proof}

\section{Reduction to the issue's statement}\label{sec:reduce0}

\begin{lemma}\label{lem:reduce-to-Bw}
Fix $c\ge1$. The following are equivalent:
\begin{enumerate}
\item[(i)] $\gamma_c(F)\subseteq N_c$;
\item[(ii)] every basic commutator of weight $\ge c$ lies in $N_c$, i.e.\
$B_{\ge c}\subseteq N_c$.
\end{enumerate}
\end{lemma}

\begin{proof}
By Proposition~\ref{prop:hall}, $B_{\ge c}$ generates $\gamma_c(F)$ as a subgroup;
since $\gamma_c(F)$ is normal, $\gamma_c(F)=\langle B_{\ge c}\rangle^{F}$. A normal
subgroup is contained in the normal subgroup $N_c$ if and only if each of its
normal generators lies in $N_c$. Hence
$\gamma_c(F)=\langle B_{\ge c}\rangle^{F}\subseteq N_c$ iff $B_{\ge c}\subseteq
N_c$.
\end{proof}

\noindent
The weight-$c$ basic commutators lie in $N_c$ trivially (they are its generators),
so the content is whether the basic commutators of weight $w>c$ do too; i.e.\
Lemma~\ref{lem:reduce-to-Bw} says $\gamma_c\subseteq N_c$ is equivalent to
$B_{>c}\subseteq N_c$. This is exactly the nontrivial point flagged in the issue:
it is \emph{not} termwise obvious, because a basic commutator $[u,v]$ of weight
$w>c$ may be built from $u,v$ of weight $<c$.

\section{Part (II): the terminating mod-$\gamma_m$ descent}\label{sec:moddescent}

\begin{lemma}[One-step descent]\label{lem:onestep}
Let $w$ be an integer with $w>c$, and suppose
\begin{equation}\label{eq:hyp-w-1}
\gamma_{w-1}(F)\subseteq N_c\,\gamma_{w}(F).
\end{equation}
Then $\gamma_{w}(F)\subseteq N_c\,\gamma_{w+1}(F)$.
\end{lemma}

\begin{proof}
By Proposition~\ref{prop:hall}, $\gamma_{w-1}(F)$ is generated as a subgroup by
$B_{\ge w-1}$; hence, being normal, $\gamma_{w-1}(F)=\langle
B_{\ge w-1}\rangle^{F}$. Apply the engine (Lemma~\ref{lem:engine}) with
$S=B_{\ge w-1}$ and $H=\gamma_{w-1}(F)$:
\[
\gamma_w(F)=[\gamma_{w-1}(F),F]
=\big\langle\,[a,x] : a\in B_{\ge w-1},\ x\in X^{\pm}\,\big\rangle^{F}.
\]
It suffices to show that for every $a\in B_{\ge w-1}$ and $x\in X^{\pm}$ the
generator $[a,x]$ lies in $N_c\,\gamma_{w+1}(F)$; since $N_c\,\gamma_{w+1}(F)$ is a
normal subgroup (a product of two normal subgroups), it then contains all
conjugates of these generators, hence all of $\gamma_w(F)$.

Fix $a\in B_{\ge w-1}$, so $a\in\gamma_{w-1}(F)$. By hypothesis
\eqref{eq:hyp-w-1} we may write $a=n\,t$ with $n\in N_c$ and $t\in\gamma_w(F)$.
By \eqref{eq:id1}, $[a,x]=[nt,x]=[n,x]^{t}\,[t,x]$. Now $[n,x]\in N_c$ by
Lemma~\ref{lem:caseA}, hence $[n,x]^{t}\in N_c$ by normality of $N_c$; and
$[t,x]\in[\gamma_w(F),F]=\gamma_{w+1}(F)$. Therefore $[a,x]\in
N_c\,\gamma_{w+1}(F)$, as required.
\end{proof}

\begin{proposition}[Mod-$\gamma_m$ form]\label{prop:moddescent}
For every integer $w\ge c$, $\ \gamma_w(F)\subseteq N_c\,\gamma_{w+1}(F)$.
Consequently $\gamma_c(F)\subseteq N_c\,\gamma_m(F)$ for every $m\ge c$.
\end{proposition}

\begin{proof}
Induction on $w\ge c$. \emph{Base $w=c$:} by \eqref{eq:hall-decomp},
$\gamma_c(F)=\langle B_c\rangle\cdot\gamma_{c+1}(F)$, and $\langle B_c\rangle
\subseteq N_c$, so $\gamma_c(F)\subseteq N_c\,\gamma_{c+1}(F)$. \emph{Step:} for
$w>c$, the case $w-1$ is $\gamma_{w-1}(F)\subseteq N_c\,\gamma_{w}(F)$, i.e.\
hypothesis \eqref{eq:hyp-w-1}; Lemma~\ref{lem:onestep} gives $\gamma_w(F)\subseteq
N_c\,\gamma_{w+1}(F)$. For the last assertion, iterate from $w=c$ to $w=m-1$, using
$N_c\cdot(N_c\,\gamma_{w+1})=N_c\,\gamma_{w+1}$ at each step (as $N_c$ is a
subgroup): $\gamma_c\subseteq N_c\gamma_{c+1}\subseteq\cdots\subseteq N_c\gamma_m$.
\end{proof}

\section{Part (III): the residue and why it needs a genuine input}\label{sec:reduce}

\subsection{Equivalent forms of the residue}

\begin{theorem}\label{thm:equiv}
Fix $c\ge1$ and let $q\colon F\twoheadrightarrow Q_c:=F/N_c$ be the quotient map.
The following are equivalent:
\begin{enumerate}
\item[(a)] $\gamma_c(F)=N_c$;
\item[(b)] $\bigcap_{m\ge c} N_c\,\gamma_m(F)=N_c$;
\item[(c)] $Q_c$ is residually nilpotent, i.e.\ $\bigcap_{k\ge1}\gamma_k(Q_c)=\{1\}$;
\item[(d)] $B_{>c}\subseteq N_c$.
\end{enumerate}
\end{theorem}

\begin{proof}
Since $q$ is surjective, $q(\gamma_k(F))=\gamma_k(Q_c)$ for all $k$ (a surjective
homomorphism carries $\gamma_k$ onto $\gamma_k$, by induction on $k$ using that it
carries generating commutators onto generating commutators); thus
$\gamma_k(Q_c)=\gamma_k(F)\,N_c/N_c$, and
$q^{-1}(\gamma_k(Q_c))=N_c\,\gamma_k(F)$.

\emph{(a)$\Leftrightarrow$(d):} this is Lemma~\ref{lem:reduce-to-Bw} together with
Lemma~\ref{lem:easy} (which makes $\gamma_c=N_c$ equivalent to
$\gamma_c\subseteq N_c$).

\emph{(b)$\Leftrightarrow$(c):} taking preimages,
\[
\bigcap_{m\ge c} N_c\,\gamma_m(F)
=\bigcap_{m\ge c} q^{-1}(\gamma_m(Q_c))
=q^{-1}\!\Big(\bigcap_{k\ge1}\gamma_k(Q_c)\Big),
\]
using that the $\gamma_k(Q_c)$ are decreasing. Hence the left side equals
$N_c=q^{-1}(\{1\})$ iff $\bigcap_k\gamma_k(Q_c)=\{1\}$.

\emph{(a)$\Leftrightarrow$(b):} set $I=\bigcap_{m\ge c}N_c\,\gamma_m(F)$.
Proposition~\ref{prop:moddescent} gives $\gamma_c(F)\subseteq N_c\,\gamma_m(F)$ for
all $m\ge c$, so $\gamma_c(F)\subseteq I$; conversely
$N_c\,\gamma_m(F)\subseteq\gamma_c(F)$ for $m\ge c$ (as $N_c\subseteq\gamma_c$ by
Lemma~\ref{lem:easy} and $\gamma_m\subseteq\gamma_c$), so $I\subseteq\gamma_c(F)$.
Hence $I=\gamma_c(F)$ always, and (b) ($I=N_c$) is equivalent to
$\gamma_c(F)=N_c$.
\end{proof}

\begin{lemma}[$N_c$ and $\gamma_c$ have identical graded data]\label{lem:gradedNc}
For every $w\ge1$, the images of $N_c$ and of $\gamma_c(F)$ in the layer
$L_w:=\gamma_w(F)/\gamma_{w+1}(F)$ coincide. Equivalently, $N_c$ and $\gamma_c(F)$
have the same image in every nilpotent quotient $F/\gamma_m(F)$.
\end{lemma}

\begin{proof}
\emph{Case $w\ge c$.} By Proposition~\ref{prop:moddescent},
$\gamma_w(F)\subseteq N_c\,\gamma_{w+1}(F)$. Since $\gamma_{w+1}(F)\subseteq
\gamma_w(F)$, Dedekind's modular law gives
\[
\gamma_w(F)=\gamma_w(F)\cap\big(N_c\,\gamma_{w+1}(F)\big)
=\big(N_c\cap\gamma_w(F)\big)\,\gamma_{w+1}(F).
\]
Passing to $L_w=\gamma_w(F)/\gamma_{w+1}(F)$, the image of $N_c$ (i.e.\ of
$N_c\cap\gamma_w(F)$) is all of $L_w$. The image of $\gamma_c(F)$ is also all of
$L_w$, since $\gamma_c(F)\supseteq\gamma_w(F)$ already surjects onto $L_w$. Hence
the two images coincide.

\emph{Case $1\le w<c$.} Then $w+1\le c$, so $N_c\subseteq\gamma_c(F)\subseteq
\gamma_{w+1}(F)$ and likewise $\gamma_c(F)\subseteq\gamma_{w+1}(F)$. Thus both
$N_c$ and $\gamma_c(F)$ have trivial image in $L_w=\gamma_w(F)/\gamma_{w+1}(F)$,
and again the two images coincide.

The reformulation in terms of nilpotent quotients $F/\gamma_m(F)$ follows because
$F/\gamma_m(F)$ is determined by the layers $L_1,\dots,L_{m-1}$ and the images of
a subgroup in those layers.
\end{proof}

\begin{example}\label{ex:Z}
Equal graded data does not force equality of subgroups, even in the presence of
residual nilpotence. Let $G=\mathbb Z$ with the central filtration
$F_w=2^{\,w-1}\mathbb Z$. Then $[F_i,F_j]=\{0\}\subseteq F_{i+j}$,
$\bigcap_w F_w=\{0\}$, and each layer $F_w/F_{w+1}\cong\mathbb Z/2$. Take
$B=\mathbb Z$ and $A=3\mathbb Z$. For every $w$,
$3\mathbb Z\cap2^{\,w-1}\mathbb Z=3\cdot2^{\,w-1}\mathbb Z$, whose image in
$F_w/F_{w+1}\cong\mathbb Z/2$ is generated by the odd multiple $3$ of
$2^{\,w-1}$, hence is all of $\mathbb Z/2$. Thus
$\mathrm{gr}_w(A)=\mathrm{gr}_w(B)=\mathbb Z/2$ for all $w$, yet
$A=3\mathbb Z\ne\mathbb Z=B$.
\end{example}

\begin{remark}[Why graded data and residual nilpotence of $F$ are insufficient]
\label{rem:graded-insufficient}
Lemma~\ref{lem:gradedNc} shows $N_c$ and $\gamma_c(F)$ are indistinguishable in
every layer $L_w$, equivalently in every nilpotent quotient $F/\gamma_m(F)$. Any
statement about the system $\{F/\gamma_m(F)\}_m$ -- about the associated graded
ring, about the Magnus power-series filtration (whose degree-$w$ part computes
$L_w$), or about residual nilpotence of $F$ itself -- is therefore \emph{blind} to
the distinction between $N_c$ and $\gamma_c(F)$. By Theorem~\ref{thm:equiv} the
missing equality is residual nilpotence of $F/N_c$, \emph{not} of $F$; and
Example~\ref{ex:Z} exhibits a residually nilpotent filtered group with free
abelian layers in which two nested subgroups with identical layer images are
nonetheless unequal. This rules out, as a matter of logic, deducing
$N_c=\gamma_c(F)$ from layer data together with residual nilpotence of the ambient
group. A genuine input about $F$ beyond its graded shadow is necessary; it is
supplied in \S\ref{sec:final} by the collection process
(Theorem~\ref{thm:hall-full}).
\end{remark}

\section{Part (IV): discharging the residue}\label{sec:final}

By Theorem~\ref{thm:equiv} the proof is complete once we establish, for every
$c\ge1$, the equivalent statements
\[
\gamma_c(F)\subseteq N_c
\quad\Longleftrightarrow\quad
B_{>c}\subseteq N_c
\quad\Longleftrightarrow\quad
F/N_c \text{ is residually nilpotent.}
\]
Remark~\ref{rem:graded-insufficient} establishes that this residue is \emph{not}
deducible from the graded Hall theorem (Proposition~\ref{prop:hall}), from the
associated graded Lie ring, or from residual nilpotence of $F$ alone
(Theorem~\ref{thm:magnus}): all of these are blind to the distinction between
$N_c$ and $\gamma_c(F)$ (Lemma~\ref{lem:gradedNc}), and
Example~\ref{ex:Z} exhibits a residually nilpotent filtered group with free
abelian layers in which two nested subgroups with identical layer data are
nevertheless unequal. A genuinely stronger input about $F$ is therefore logically
necessary. We supply it as a single, precisely stated, classical theorem, verify
its hypotheses, and give the complete bridge to Theorem~\ref{thm:main}. We are
explicit (Remark~\ref{rem:honest}) about the fact that the substantive content of
the closure step is \emph{imported}, not reproved from the graded data.

\subsection{The substantive input}

\begin{theorem}[Normal-closure form of the Hall basis theorem]\label{thm:hall-full}
Let $F$ be free of finite rank $r$ with free basis $X=\{x_1,\dots,x_r\}$, and use
the Hall ordering and weighting of basic commutators of the problem statement. Then
for every $w\ge1$,
\begin{equation}\label{eq:hall-claim}
\gamma_w(F)=\langle B_w\rangle^{F}=N_w .
\end{equation}
Equivalently, every basic commutator of weight $m\ge w$ lies in $N_w$, i.e.\
$B_{\ge w}\subseteq N_w$.
\end{theorem}

\noindent
\emph{Reference and precise content.}
This is Marshall Hall, \emph{The Theory of Groups} (Macmillan, 1959),
Chapter~11: the Collecting Process (\S11.1, Theorem~11.1.4) together with the
Basis Theorem (\S11.2, Theorem~11.2.4) and its discussion of normal generation;
equivalently W.\ Magnus, A.\ Karrass, D.\ Solitar, \emph{Combinatorial Group
Theory} (2nd ed., Dover, 1976), \S5.3--5.7, Theorem~5.13A with its corollaries.
The statement \eqref{eq:hall-claim} is the on-the-nose (normal-closure) refinement
of Proposition~\ref{prop:hall}: Proposition~\ref{prop:hall} asserts that $B_{\ge
w}$ \emph{generates $\gamma_w(F)$ as a subgroup} and that the layers
$\gamma_w/\gamma_{w+1}$ are free abelian on $B_w$, whereas
Theorem~\ref{thm:hall-full} asserts in addition that the higher-weight generators
$B_{>w}$ are already contained in the normal closure $N_w=\langle B_w\rangle^F$.

\smallskip
\emph{Verification of hypotheses.} The hypotheses of
Theorem~\ref{thm:hall-full} are: $F$ free of finite rank $r$; $X$ the given free
basis; and $B_w$ the set of basic commutators of weight exactly $w$ in the Hall
ordering, with the weight function and ordering defined in the problem statement.
These are verbatim the hypotheses of the present problem (\S\ref{sec:input}). In
particular the ordering used in M.\ Hall's collecting process is exactly the
weight-respecting Hall ordering specified here, so the theorem applies with no
adjustment.

\smallskip
\emph{Source of the input, and why it is not vacuously the present theorem.}
The proof of Theorem~\ref{thm:hall-full} is the \emph{collection process}: a
terminating rewriting procedure on words in $X^{\pm}$ which, for each $g\in F$ and
each $n\ge1$, rewrites $g$ into its \emph{collected form}
\begin{equation}\label{eq:collected}
g\equiv b_{1}^{\,a_{1}}\,b_{2}^{\,a_{2}}\cdots b_{k(n)}^{\,a_{k(n)}}
\pmod{\gamma_{n+1}(F)},
\end{equation}
where $b_1<b_2<\cdots$ is the increasing Hall enumeration of all basic
commutators, $b_1,\dots,b_{k(n)}$ are those of weight $\le n$, the exponents
$a_i\in\mathbb Z$ are uniquely determined by $g$, and---this is the genuinely
non-graded output of the algorithm---each $a_i$ is independent of $n$ (enlarging
$n$ only appends factors of larger weight). The procedure is defined and proved
correct purely by repeated application of the commutator identities
\eqref{eq:id1}--\eqref{eq:id2} and well-founded induction on the collection
ordering; it refers neither to $N_w$ nor to \eqref{eq:hall-claim}. From this
algorithm \eqref{eq:hall-claim} is obtained because the rewriting that produces a
weight-$m$ collector ($m>w$) from a word lying in $\gamma_w(F)$ does so by
commuting a weight-$w$ basic commutator past generators, i.e.\ it exhibits each
weight-$m$ collector as a product of conjugates of weight-$w$ basic commutators.
We do not reproduce the (lengthy) collection algorithm here; we cite it as the
classical theorem it is and use only its stated conclusion \eqref{eq:hall-claim}.

\begin{lemma}[The closure inclusion]\label{lem:hall-residual}
For every $c\ge1$, $\ B_{>c}\subseteq N_c$, and hence $\gamma_c(F)\subseteq N_c$.
\end{lemma}

\begin{proof}
Apply Theorem~\ref{thm:hall-full} with $w=c$: every basic commutator of weight
$m\ge c$ lies in $N_c$, i.e.\ $B_{\ge c}\subseteq N_c$; in particular
$B_{>c}\subseteq N_c$. Lemma~\ref{lem:reduce-to-Bw} then gives
$\gamma_c(F)\subseteq N_c$. The hypotheses of Theorem~\ref{thm:hall-full} were
verified above.
\end{proof}

\subsection{The bridge and conclusion}\label{ssec:bridge}

\begin{corollary}[Residual nilpotence of $F/N_c$]\label{cor:resnil}
For every $c\ge1$, $\ \bigcap_{m\ge c}N_c\,\gamma_m(F)=N_c$, and $F/N_c$ is
residually nilpotent.
\end{corollary}

\begin{proof}
By Lemma~\ref{lem:hall-residual}, $\gamma_c(F)\subseteq N_c$; with
Lemma~\ref{lem:easy} this gives $N_c=\gamma_c(F)$. Hence for $m\ge c$,
$N_c\,\gamma_m(F)=\gamma_c(F)\,\gamma_m(F)=\gamma_c(F)=N_c$ (using
$\gamma_m\subseteq\gamma_c$), so $\bigcap_{m\ge c}N_c\,\gamma_m(F)=N_c$. By
Theorem~\ref{thm:equiv}, (b)$\Rightarrow$(c), $F/N_c$ is residually nilpotent.
\end{proof}

\begin{proof}[Proof of Theorem~\ref{thm:main}]
The inclusion $N_c\subseteq\gamma_c(F)$ is Lemma~\ref{lem:easy}. The reverse
inclusion $\gamma_c(F)\subseteq N_c$ is Lemma~\ref{lem:hall-residual}. Hence
$\gamma_c(F)=N_c$ for every $c\ge1$, and the answer to the problem is \emph{yes}.
\end{proof}

\begin{remark}[Honest account of what is imported and what is proved]
\label{rem:honest}
We state precisely the logical economy of the proof, without overclaiming.

\emph{(1) Proved here in full}, from Proposition~\ref{prop:hall} (graded Hall),
Theorem~\ref{thm:magnus} (Magnus), and elementary commutator calculus:
the easy inclusion $N_c\subseteq\gamma_c(F)$ (Lemma~\ref{lem:easy}); the
commutator engine (Lemma~\ref{lem:engine}); the mod-$\gamma_m$ descent
$\gamma_c\subseteq N_c\gamma_m$ for all $m$ (Proposition~\ref{prop:moddescent});
the equivalence of the goal with $B_{>c}\subseteq N_c$, with
$\bigcap_m N_c\gamma_m=N_c$, and with residual nilpotence of $F/N_c$
(Theorem~\ref{thm:equiv}); the fact that $N_c$ and $\gamma_c(F)$ have identical
graded data (Lemma~\ref{lem:gradedNc}); and the \emph{necessity} statement
(Remark~\ref{rem:graded-insufficient}, Example~\ref{ex:Z}) that the residue is not
obtainable from graded data nor from residual nilpotence of $F$.

\emph{(2) Imported as a classical theorem.} The single substantive input is
Theorem~\ref{thm:hall-full}, the normal-closure form of the Hall basis theorem.
We do not reprove it; we cite it (M.\ Hall, \emph{Theory of Groups}, \S11.1--11.2;
MKS \S5.3--5.7), verify its hypotheses are exactly those of the problem, and use
its conclusion \eqref{eq:hall-claim} through Lemma~\ref{lem:hall-residual}.

\emph{(3) Honest assessment of the import.} We make no claim that
Theorem~\ref{thm:hall-full} is strictly weaker than, or logically independent of,
the present Theorem~\ref{thm:main}: by Theorem~\ref{thm:equiv} the conclusion
$\gamma_c(F)\subseteq N_c$ is \emph{equivalent} to residual nilpotence of $F/N_c$,
and the cited theorem supplies exactly this. Thus the substantive heart of the
closure step is \emph{imported}, and the value of \S\S\ref{sec:engine}--\ref{sec:reduce}
is to (i) reduce the problem to a single, sharply identified residue, (ii) prove
rigorously that this residue cannot be discharged by graded methods or by residual
nilpotence of $F$, thereby pinpointing the precise additional strength required,
and (iii) connect that residue cleanly to the classical theorem that provides it.
The one genuinely non-graded fact the collection process contributes beyond
Proposition~\ref{prop:hall} is the consistency of the collected exponents across
all truncations in \eqref{eq:collected}; this is information unavailable from any
single layer $\gamma_w/\gamma_{w+1}$ or finite nilpotent quotient
$F/\gamma_m(F)$, and it is precisely what upgrades
$\gamma_c\subseteq\bigcap_m N_c\gamma_m$ (Proposition~\ref{prop:moddescent}) to
$\gamma_c\subseteq N_c$.
\end{remark}

\end{document}
